\chapter{高等数学：保号性}

\section{数列与函数的最终保号}

\begin{problemblock}
\textbf{例1.} 设 \(x_1>0\)，数列 \(\{x_n\}\) 满足
\[
\lim_{n\to\infty}\frac{x_{n+1}}{x_n}=\frac12,
\]
证明：
\[
\lim_{n\to\infty}x_n=0.
\]
\end{problemblock}
\begin{solutionblock}
\analysis{比值极限小于 \(1\)，就能推出从某一项开始按几何级数衰减。}
由
\[
\lim_{n\to\infty}\frac{x_{n+1}}{x_n}=\frac12
\]
可知，取 \(\frac34\) 作为上界，存在 \(N\)，使得当 \(n\ge N\) 时
\[
\left|\frac{x_{n+1}}{x_n}\right|<\frac34.
\]
于是
\[
|x_{n+1}|<\frac34|x_n|,\qquad n\ge N.
\]
递推得
\[
|x_n|\le |x_N|\left(\frac34\right)^{n-N},\qquad n\ge N.
\]
由于
\[
\left(\frac34\right)^{n-N}\to0,
\]
由夹逼定理得
\[
\lim_{n\to\infty}x_n=0.
\]
\answer{已证 \(\displaystyle \lim_{n\to\infty}x_n=0\)}
\examnote{比值极限为 \(q\) 且 \(|q|<1\)，通常直接取一个介于 \(|q|\) 与 \(1\) 之间的数做最终控制。}
\end{solutionblock}

\begin{problemblock}
\textbf{例2.} 设数列 \(\{x_n\},\{y_n\}\) 满足
\[
x_1=y_1=\frac12,\qquad x_{n+1}=\sin x_n,\qquad y_{n+1}=y_n^2.
\]
当 \(n\to\infty\) 时，证明：\(y_n\) 是 \(x_n\) 的高阶无穷小。
\end{problemblock}
\begin{solutionblock}
\analysis{这里要证 \(y_n=o(x_n)\)。其中 \(y_n\) 是双指数衰减，而 \(x_{n+1}/x_n\to1\)，所以 \(x_n\) 衰减得远慢于 \(y_n\)。}
先看 \(x_n\)。由于 \(0<x_1<1\)，且当 \(0<x<1\) 时
\[
0<\sin x<x,
\]
所以 \(x_n>0\)、单调递减，且存在极限。设极限为 \(l\ge0\)，由递推式得
\[
l=\sin l,
\]
故 \(l=0\)。

又
\[
\lim_{n\to\infty}\frac{x_{n+1}}{x_n}
=\lim_{n\to\infty}\frac{\sin x_n}{x_n}=1.
\]
因此存在 \(N\)，使得 \(n\ge N\) 时
\[
x_{n+1}>\frac12x_n.
\]
于是对 \(n\ge N\)，有
\[
x_n\ge x_N\left(\frac12\right)^{n-N}.
\]
另一方面
\[
y_n=\left(\frac12\right)^{2^{n-1}}.
\]
故
\[
0\le \frac{y_n}{x_n}
\le \frac{2^N}{x_N}\left(\frac12\right)^{2^{n-1}-n}\to0.
\]
所以
\[
y_n=o(x_n).
\]
\answer{\(\displaystyle y_n=o(x_n)\)}
\examnote{\(\sin x\sim x\) 给的是普通几何级别衰减，而 \(y_{n+1}=y_n^2\) 给的是双指数衰减，后者远快于前者。}
\end{solutionblock}

\begin{problemblock}
\textbf{例3.} 已知数列 \(\{x_n\},\{y_n\}\) 满足
\[
x_1=y_1=\frac12,\qquad x_{n+1}=1-\cos x_n,\qquad y_{n+1}=y_n^2
\quad (n=1,2,\cdots).
\]
则当 \(n\to\infty\) 时，（\quad）

A. \(\displaystyle \lim_{n\to\infty}\frac{x_n-y_n}{y_n}=0\)
\quad
B. \(\displaystyle \lim_{n\to\infty}\frac{x_n+y_n}{y_n}=1\)

C. \(\displaystyle \lim_{n\to\infty}\frac{x_n-y_n}{x_n}=0\)
\quad
D. \(\displaystyle \lim_{n\to\infty}\frac{x_n+y_n}{x_n}=1\)
\end{problemblock}
\begin{solutionblock}
\analysis{比较 \(x_n\) 与 \(y_n\) 的大小。因为 \(1-\cos x\le x^2/2\)，所以 \(x_n\) 比 \(y_n\) 还要小一个迭代量级。}
令
\[
r_n=\frac{x_n}{y_n}.
\]
因为 \(x_1=y_1=1/2\)，所以 \(r_1=1\)。又
\[
1-\cos x\le \frac{x^2}{2},
\]
故
\[
r_{n+1}
=\frac{x_{n+1}}{y_{n+1}}
=\frac{1-\cos x_n}{y_n^2}
\le \frac{x_n^2}{2y_n^2}
=\frac12 r_n^2.
\]
由 \(r_1=1\) 得
\[
r_2\le\frac12,\qquad r_3\le\frac12\left(\frac12\right)^2,
\]
递推可知
\[
r_n\to0.
\]
即
\[
x_n=o(y_n).
\]
于是
\[
\frac{x_n+y_n}{y_n}=1+\frac{x_n}{y_n}\to1.
\]
\answer{B}
\examnote{遇到 \(1-\cos x_n\) 的递推，核心等价是 \(1-\cos x\sim x^2/2\)，但证明不等式时用 \(1-\cos x\le x^2/2\) 更稳。}
\end{solutionblock}

\begin{problemblock}
\textbf{例4.} 定义在实数域上的连续函数 \(f(x)\) 满足
\[
f(2-x)=f(x),
\]
且 \(x_0=2\) 为 \(f(x)\) 在 \((1,+\infty)\) 上的唯一零点，\(f'(0)<0\)。
设
\[
F(x)=\int_0^x f(t)\,dt,
\]
则 \(F(x)\) 的严格单调递减区间为（\quad）

A. \((0,2)\)
\quad B. \((2,+\infty)\)
\quad C. \((-\infty,0)\)
\quad D. 无法确定
\end{problemblock}
\begin{solutionblock}
\analysis{\(F'(x)=f(x)\)，所以要找 \(f(x)<0\) 的区间。题中的对称性 \(f(2-x)=f(x)\) 是关于 \(x=1\) 对称。}
由 \(x_0=2\) 为零点，得
\[
f(2)=0.
\]
又
\[
f(2-x)=f(x),
\]
令 \(x=0\)，得
\[
f(2)=f(0),
\]
所以
\[
f(0)=0.
\]
由于
\[
f'(0)<0,
\]
故在 \(0\) 的右邻域内 \(f(x)<0\)，在 \(0\) 的左邻域内 \(f(x)>0\)。

对称性给出
\[
f(2-h)=f(h),\qquad f(2+h)=f(-h).
\]
因此在 \(2\) 的左邻域内 \(f<0\)，在 \(2\) 的右邻域内 \(f>0\)。
又 \(2\) 是 \((1,+\infty)\) 上唯一零点，所以
\[
f(x)<0\quad (1<x<2),\qquad f(x)>0\quad (x>2).
\]
再由对称性可得
\[
f(x)<0\quad (0<x<1),\qquad f(x)>0\quad (x<0).
\]
于是
\[
f(x)<0\quad (0<x<2).
\]
因为
\[
F'(x)=f(x),
\]
所以 \(F\) 的严格单调递减区间为
\[
(0,2).
\]
\answer{A}
\examnote{积分函数单调性先看导函数符号；本题真正的工作是借对称性和唯一零点确定 \(f\) 的符号。}
\end{solutionblock}

\begin{problemblock}
\textbf{例5.} 设函数 \(f(x)\) 在 \([0,+\infty)\) 上具有一阶连续导数，
\[
f(0)=0,\qquad \lim_{x\to+\infty}\bigl[f(x)+f'(x)\bigr]=a>0.
\]
证明：存在 \(x_0>0\)，使得当 \(x>x_0\) 时
\[
f(x)>\frac a2.
\]
\end{problemblock}
\begin{solutionblock}
\analysis{把 \(f+f'\) 看成 \((e^xf)'\) 的核心部分。极限为正意味着从某处开始 \((e^xf)'\) 至少像正的 \(e^x\) 那样增长。}
由
\[
\lim_{x\to+\infty}\bigl[f(x)+f'(x)\bigr]=a>0,
\]
取 \(\frac{3a}{4}\)，存在 \(X>0\)，当 \(x>X\) 时
\[
f(x)+f'(x)>\frac{3a}{4}.
\]
注意
\[
(e^xf(x))'=e^x(f(x)+f'(x)).
\]
故当 \(x>X\) 时，
\[
(e^xf(x))'>\frac{3a}{4}e^x.
\]
从 \(X\) 到 \(x\) 积分，得
\[
e^xf(x)-e^Xf(X)
>\frac{3a}{4}(e^x-e^X).
\]
两边除以 \(e^x\)，
\[
f(x)>\frac{3a}{4}
+e^{X-x}\left(f(X)-\frac{3a}{4}\right).
\]
令 \(x\to+\infty\)，右端趋于 \(3a/4\)。因此存在 \(x_0>X\)，当 \(x>x_0\) 时
\[
f(x)>\frac a2.
\]
\answer{已证存在 \(x_0>0\)，使 \(x>x_0\) 时 \(f(x)>a/2\)}
\examnote{看到 \(f+f'\) 的最终符号，优先乘以积分因子 \(e^x\)。}
\end{solutionblock}

\begin{problemblock}
\textbf{例6.} 已知 \(f(x)\) 为 \((0,+\infty)\) 上连续的有界函数，且满足
\[
\lim_{x\to0^+}\frac{f(x)}{x-\tan x}=6.
\]
证明：方程
\[
f(x)+x^3=0
\]
在 \((0,+\infty)\) 上必有一实根。
\end{problemblock}
\begin{solutionblock}
\analysis{先用 \(x-\tan x\sim -x^3/3\) 判断 \(f(x)+x^3\) 在 \(0^+\) 附近为负，再用有界性判断它在充分大时为正。}
当 \(x\to0^+\) 时，
\[
\tan x=x+\frac{x^3}{3}+o(x^3),
\]
所以
\[
x-\tan x=-\frac{x^3}{3}+o(x^3).
\]
由题设极限，
\[
f(x)\sim 6(x-\tan x)\sim -2x^3.
\]
因此
\[
f(x)+x^3\sim -x^3<0\qquad (x\to0^+).
\]
故存在 \(\delta>0\)，使得
\[
0<x<\delta
\]
时
\[
f(x)+x^3<0.
\]
又因为 \(f\) 有界，存在 \(M>0\)，使得
\[
|f(x)|\le M.
\]
当 \(x>\sqrt[3]{M}\) 时，
\[
f(x)+x^3\ge -M+x^3>0.
\]
函数 \(f(x)+x^3\) 在 \((0,+\infty)\) 上连续，故由介值定理，方程
\[
f(x)+x^3=0
\]
在 \((0,+\infty)\) 上至少有一个实根。
\answer{已证方程在 \((0,+\infty)\) 上必有实根}
\examnote{存在根问题常用“近端一个符号、远端一个符号”。这里近端靠等价无穷小，远端靠有界性。}
\end{solutionblock}

\section{导数极限、极值与拐点}

\begin{problemblock}
\textbf{例7.} \(f(x)\) 在 \([0,+\infty)\) 上连续可导，且
\[
\lim_{x\to+\infty}f'(x)
\]
存在。若 \(f(x)\) 在 \([0,+\infty)\) 上有界，证明：
\[
\lim_{x\to+\infty}f'(x)=0.
\]
\end{problemblock}
\begin{solutionblock}
\analysis{若导数极限不是 \(0\)，则导数最终保持同号且绝对值有正下界，从而函数必然单调发散，这与有界矛盾。}
设
\[
\lim_{x\to+\infty}f'(x)=A.
\]
若 \(A>0\)，则存在 \(X>0\)，使得当 \(x>X\) 时
\[
f'(x)>\frac A2.
\]
于是
\[
f(x)-f(X)=\int_X^x f'(t)\,dt>\frac A2(x-X)\to+\infty,
\]
与 \(f\) 有界矛盾。

若 \(A<0\)，同理存在 \(X>0\)，使得当 \(x>X\) 时
\[
f'(x)<\frac A2<0,
\]
从而 \(f(x)\to-\infty\)，仍与有界矛盾。

因此只能有
\[
A=0.
\]
\answer{\(\displaystyle \lim_{x\to+\infty}f'(x)=0\)}
\examnote{“有界 + 导数极限存在”常用反证：导数若最终离开 0，函数就会线性发散。}
\end{solutionblock}

\begin{problemblock}
\textbf{例8.} 设函数 \(f(x)\) 在区间 \((0,+\infty)\) 内有界且存在二阶导数，则下列说法正确的是（\quad）

A. 若 \(\displaystyle \lim_{x\to+\infty}f(x)\) 存在，则 \(\displaystyle \lim_{x\to+\infty}f''(x)\) 存在。

B. 若 \(\displaystyle \lim_{x\to+\infty}f''(x)\) 存在，则 \(\displaystyle \lim_{x\to+\infty}f(x)\) 存在。

C. 若 \(\displaystyle \lim_{x\to+\infty}f''(x)\) 存在，则 \(\displaystyle \lim_{x\to+\infty}f''(x)=0\)。

D. 若 \(\displaystyle \lim_{x\to+\infty}f'(x)=0\) 存在，则 \(\displaystyle \lim_{x\to+\infty}f''(x)=0\)。
\end{problemblock}
\begin{solutionblock}
\analysis{有界性只能排除 \(f''\) 有非零极限；不能保证 \(f\)、\(f'\)、\(f''\) 的其他极限都存在。}
设
\[
\lim_{x\to+\infty}f''(x)=A.
\]
若 \(A>0\)，则存在 \(X\)，使得 \(x>X\) 时
\[
f''(x)>\frac A2.
\]
积分一次得 \(f'(x)\) 至少线性增长，再积分一次得 \(f(x)\) 至少二次增长，这与 \(f\) 有界矛盾。

若 \(A<0\)，同理 \(f(x)\) 会向 \(-\infty\) 方向二次发散，也与有界矛盾。
故只能
\[
A=0.
\]
所以 C 正确。

其余选项不正确：例如
\[
f(x)=\sin(\ln x)
\]
有界且 \(f''(x)\to0\)，但 \(f(x)\) 无极限，排除 B；
\[
f(x)=\frac{\sin x^2}{x}
\]
在 \((0,+\infty)\) 上有界且 \(f(x)\to0\)，但 \(f''(x)\) 无极限，排除 A；再取
\[
f(x)=\int_1^x\frac{\sin t^2}{t}\,dt,
\]
则 \(f\) 有界，\(f'(x)=\sin x^2/x\to0\)，但
\[
f''(x)=2\cos x^2-\frac{\sin x^2}{x^2}
\]
无极限，排除 D。
\answer{C}
\examnote{本题只需要证明“若 \(f''\) 极限存在，则不能非零”。其他更强结论都容易被振荡函数反例推翻。}
\end{solutionblock}

\begin{problemblock}
\textbf{例9.} 已知
\[
a_n=\frac{\sin n}{n}\qquad (n=1,2,\cdots),
\]
则数列 \(\{a_n\}\)（\quad）

A. 有最大值，有最小值。
\quad B. 有最大值，没有最小值。

C. 没有最大值，有最小值。
\quad D. 没有最大值，没有最小值。
\end{problemblock}
\begin{solutionblock}
\analysis{数列 \(\sin n/n\to0\)。只要能找到正项和负项，就可把极值问题压缩到有限项中处理。}
显然
\[
\lim_{n\to\infty}\frac{\sin n}{n}=0.
\]
又
\[
a_1=\sin1>0,
\]
而例如
\[
a_4=\frac{\sin4}{4}<0.
\]
由于 \(a_n\to0\)，存在 \(N\)，使得 \(n>N\) 时
\[
|a_n|<\frac{a_1}{2},
\]
所以最大值只能在有限集合
\[
\{a_1,a_2,\cdots,a_N\}
\]
中取得。

同理，因为 \(a_4<0\)，又 \(a_n\to0\)，可取 \(N_1\)，使得 \(n>N_1\) 时
\[
a_n>\frac{a_4}{2}.
\]
注意 \(a_4/2<0\) 且 \(a_4<a_4/2\)，所以最小值也只能在有限集合
\[
\{a_1,a_2,\cdots,a_{N_1}\}
\]
中取得。

因此数列既有最大值，也有最小值。
\answer{A}
\examnote{收敛数列的最大最小不一定都存在；本题能存在，是因为极限为 \(0\)，同时已经出现正项和负项，可把极值锁定在有限前段。}
\end{solutionblock}

\begin{problemblock}
\textbf{例10.} 设函数 \(f(x)\) 连续，给出下列四个条件：
\[
\begin{array}{ll}
(1)&\displaystyle \lim_{x\to0}\frac{|f(x)|-f(0)}{x}\ \text{存在};\\[0.8em]
(2)&\displaystyle \lim_{x\to0}\frac{f(x)-|f(0)|}{x}\ \text{存在};\\[0.8em]
(3)&\displaystyle \lim_{x\to0}\frac{|f(x)|}{x}\ \text{存在};\\[0.8em]
(4)&\displaystyle \lim_{x\to0}\frac{|f(x)|-|f(0)|}{x}\ \text{存在}.
\end{array}
\]
能得到“\(f(x)\) 在 \(x=0\) 处可导”的条件个数是（\quad）

A. 1
\quad B. 2
\quad C. 3
\quad D. 4
\end{problemblock}
\begin{solutionblock}
\analysis{逐个讨论 \(f(0)\) 的符号。若 \(f(0)\ne0\)，连续性保证 \(f\) 在 0 附近不变号；若 \(f(0)=0\)，两侧极限的符号会迫使绝对值商的极限为 0。}
(1) 若 \(f(0)>0\)，则 \(0\) 附近 \(|f(x)|=f(x)\)，条件 (1) 就是
\[
\lim_{x\to0}\frac{f(x)-f(0)}x
\]
存在。若 \(f(0)<0\)，则 \(|f(x)|-f(0)\to -2f(0)\ne0\)，商不可能有有限极限。若 \(f(0)=0\)，条件 (1) 变为
\[
\lim_{x\to0}\frac{|f(x)|}{x}
\]
存在；由于 \(x>0\) 时该商非负，\(x<0\) 时该商非正，两侧极限若相等只能为 \(0\)，从而
\[
\frac{f(x)}x\to0.
\]
故 (1) 能推出可导。

(2) 若 \(f(0)\ge0\)，则 \(|f(0)|=f(0)\)，条件 (2) 正是可导定义；若 \(f(0)<0\)，分子趋于 \(2f(0)\ne0\)，极限不可能存在。因此 (2) 也能推出可导。

(3) 若
\[
\lim_{x\to0}\frac{|f(x)|}{x}
\]
存在，则首先必须有 \(f(0)=0\)。同样由左右符号可知该极限只能为 \(0\)，于是
\[
\left|\frac{f(x)}x\right|\to0,
\]
故 \(f'(0)=0\)。所以 (3) 能推出可导。

(4) 若 \(f(0)\ne0\)，则 \(f\) 在 0 附近不变号，\(|f|\) 可导等价于 \(f\) 可导；若 \(f(0)=0\)，则退化为 (3)。所以 (4) 也能推出可导。

四个条件均可推出 \(f\) 在 \(0\) 处可导。
\answer{D}
\examnote{绝对值问题的分界点是 \(f(0)=0\)。若 \(f(0)\ne0\)，连续性给出局部保号；若 \(f(0)=0\)，必须利用左右极限符号。}
\end{solutionblock}

\begin{problemblock}
\textbf{例11.} 设 \(f(x)\) 具有二阶连续导数，且
\[
f'(1)=0,\qquad
\lim_{x\to1}\frac{f''(x)}{(x-1)^2}=\frac12,
\]
则（\quad）

A. \(f(1)\) 是 \(f(x)\) 的极大值。

B. \(f(1)\) 是 \(f(x)\) 的极小值。

C. \((1,f(1))\) 是曲线 \(y=f(x)\) 的拐点。

D. \(f(1)\) 不是 \(f(x)\) 的极值，\((1,f(1))\) 也不是曲线 \(y=f(x)\) 的拐点。
\end{problemblock}
\begin{solutionblock}
\analysis{由极限可知 \(f''(x)\) 在 \(x=1\) 两侧都为正，因此 \(f'\) 左负右正，得到极小值；二阶导不变号，所以不是拐点。}
由
\[
\lim_{x\to1}\frac{f''(x)}{(x-1)^2}=\frac12>0
\]
可知，存在 \(\delta>0\)，当 \(0<|x-1|<\delta\) 时
\[
f''(x)>0.
\]
因为 \(f'(1)=0\)，所以当 \(x>1\) 且接近 \(1\) 时，
\[
f'(x)=\int_1^x f''(t)\,dt>0;
\]
当 \(x<1\) 且接近 \(1\) 时，
\[
f'(x)=\int_1^x f''(t)\,dt<0.
\]
因此 \(f\) 在 \(1\) 左侧递减、右侧递增，故 \(f(1)\) 是极小值。

同时 \(f''(x)\) 在 \(1\) 两侧同为正，不发生凹凸性改变，所以 \((1,f(1))\) 不是拐点。
\answer{B}
\examnote{判断极值看 \(f'\) 变号；判断拐点看 \(f''\) 是否变号。}
\end{solutionblock}

\begin{problemblock}
\textbf{例12.} 设函数 \(f(x)\) 二阶可导，满足
\[
f'(a)=0,\qquad
\lim_{x\to a}\frac{f''(x)}{x-a}=1,
\]
则（\quad）

A. \(x=a\) 是 \(f(x)\) 的极小值点。

B. \(x=a\) 是 \(f(x)\) 的极大值点。

C. \((a,f(a))\) 是曲线 \(y=f(x)\) 的拐点。

D. \(x=a\) 不是 \(f(x)\) 的极值点，\((a,f(a))\) 也不是曲线 \(y=f(x)\) 的拐点。
\end{problemblock}
\begin{solutionblock}
\analysis{这里 \(f''(x)\sim x-a\)，所以二阶导在 \(a\) 两侧变号；但 \(f'(x)\sim (x-a)^2/2\)，一阶导两侧同号。}
由
\[
\frac{f''(x)}{x-a}\to1
\]
可知，\(x\) 充分接近 \(a\) 时，
\[
f''(x)
\]
与 \(x-a\) 同号。因此 \(f''\) 在 \(a\) 左侧为负、右侧为正，曲线凹凸性发生改变，所以
\[
(a,f(a))
\]
是拐点。

再看极值。由 \(f'(a)=0\)，且
\[
f'(x)-f'(a)=\int_a^x f''(t)\,dt.
\]
由于 \(f''(t)\sim t-a\)，故
\[
f'(x)\sim \int_a^x (t-a)\,dt=\frac{(x-a)^2}{2}>0
\]
在 \(a\) 两侧同号。故 \(x=a\) 不是极值点。
\answer{C}
\examnote{\(f''\) 变号给拐点，不一定给极值；极值还要看 \(f'\) 是否变号。}
\end{solutionblock}

\begin{problemblock}
\textbf{例13.} 设 \(f(x)\) 在 \([a,b]\) 上可导，且在 \(x=a\) 处取最小值，在 \(x=b\) 处取最大值，则（\quad）

A. \(f'_+(a)\le0,\ f'_-(b)\le0\)
\quad B. \(f'_+(a)\le0,\ f'_-(b)\ge0\)

C. \(f'_+(a)>0,\ f'_-(b)>0\)
\quad D. \(f'_+(a)\ge0,\ f'_-(b)\ge0\)
\end{problemblock}
\begin{solutionblock}
\analysis{端点极值只看单侧导数。左端点取最小值，向右的函数增量非负；右端点取最大值，从左逼近时分子和分母都非正。}
因为 \(x=a\) 处取最小值，所以当 \(x>a\) 时
\[
f(x)-f(a)\ge0,\qquad x-a>0.
\]
因此
\[
f'_+(a)=\lim_{x\to a^+}\frac{f(x)-f(a)}{x-a}\ge0.
\]
因为 \(x=b\) 处取最大值，所以当 \(x<b\) 时
\[
f(x)-f(b)\le0,\qquad x-b<0.
\]
故
\[
\frac{f(x)-f(b)}{x-b}\ge0,
\]
从而
\[
f'_-(b)\ge0.
\]
\answer{D}
\examnote{端点处没有双侧邻域，不能套内部极值 \(f'=0\)，只能看单侧差商符号。}
\end{solutionblock}

\begin{problemblock}
\textbf{例14.} 设
\[
\lim_{x\to a}\frac{f(x)-f(a)}{(x-a)^n}=-1,
\]
其中 \(n\in\mathbb N^+\)，\(n>1\)。则在 \(x=a\) 处（\quad）

A. \(f(x)\) 的导数存在，且 \(f'(a)\ne0\)。

B. \(f(x)\) 取得极大值。

C. \(f(x)\) 取得极小值。

D. \(f(x)\) 是否取得极值与 \(n\) 的取值有关。
\end{problemblock}
\begin{solutionblock}
\analysis{由极限可知 \(f(x)-f(a)\sim-(x-a)^n\)。当 \(n\) 为偶数时两侧同负，有极大值；当 \(n\) 为奇数时两侧异号，无极值。}
由题设，
\[
f(x)-f(a)\sim -(x-a)^n.
\]
若 \(n\) 为偶数，则
\[
(x-a)^n>0\quad (x\ne a),
\]
所以
\[
f(x)-f(a)<0
\]
在 \(a\) 附近两侧成立，故 \(f(a)\) 为极大值。

若 \(n\) 为奇数，则 \(x<a\) 时 \((x-a)^n<0\)，从而
\[
f(x)-f(a)>0;
\]
而 \(x>a\) 时
\[
f(x)-f(a)<0.
\]
函数值穿过 \(f(a)\)，不取得极值。

因此是否取得极值与 \(n\) 的奇偶有关。
\answer{D}
\examnote{高阶主部判断极值时，偶次看系数正负，奇次通常不是极值。}
\end{solutionblock}

\begin{problemblock}
\textbf{例15.} 设 \(f(x)\) 在 \(x=0\) 的邻域内有定义，且
\[
\lim_{x\to0}\frac{f(x)}{x-\ln(1+x)}=1,
\]
则（\quad）

A. \(f(0)=0\)
\quad B. \(f'(0)=0\)
\quad C. \(x=0\) 是 \(f(x)\) 的极小值点
\quad D. 以上说法均不正确。
\end{problemblock}
\begin{solutionblock}
\analysis{极限只控制 \(x\ne0\) 时的函数值，并没有给出 \(f\) 在 0 处连续。不能把去心极限直接当成 \(f(0)\)。}
当 \(x\to0\) 时，
\[
\ln(1+x)=x-\frac{x^2}{2}+o(x^2),
\]
所以
\[
x-\ln(1+x)=\frac{x^2}{2}+o(x^2)>0
\]
在 \(0\) 的去心邻域内成立。因此题设说明
\[
f(x)\sim \frac{x^2}{2}\qquad (x\to0,\ x\ne0).
\]
这只能推出去心极限
\[
\lim_{x\to0,\ x\ne0}f(x)=0,
\]
不能推出 \(f(0)=0\)，因为题目没有给 \(f\) 在 \(0\) 处连续。

例如可令 \(x\ne0\) 时
\[
f(x)=x-\ln(1+x),
\]
但任意指定 \(f(0)=1\)。此时题设极限仍为 \(1\)，但 A、B、C 均不成立。
\answer{D}
\examnote{涉及极值、导数、函数值时，必须区分“去心极限”与“点处函数值”。}
\end{solutionblock}

\begin{problemblock}
\textbf{例16.} 设函数 \(f(x)\) 在 \(x=0\) 处连续，在 \(x=0\) 的某个去心邻域内可导，且
\[
\lim_{x\to0}\frac{f(x)}x
\]
存在，则下列命题中，错误的是（\quad）

A. \(f(0)=0\)。

B. \(f'(0)\) 存在。

C. \(\displaystyle \lim_{x\to0}f'(x)=0\)。

D. 在 \(x=0\) 的某个去心邻域内，存在常数 \(C>0\)，使得
\[
|f(x)|<C|x|.
\]
\end{problemblock}
\begin{solutionblock}
\analysis{极限 \(f(x)/x\) 存在，加上 \(f\) 在 0 连续，可推出 \(f(0)=0\) 和 \(f'(0)\) 存在；但不能推出去心导数极限为 0。}
设
\[
\lim_{x\to0}\frac{f(x)}x=L.
\]
因为 \(x\to0\) 时 \(x\to0\)，所以
\[
f(x)=x\cdot\frac{f(x)}x\to0.
\]
又 \(f\) 在 \(0\) 处连续，故
\[
f(0)=0,
\]
A 正确。

于是
\[
f'(0)=\lim_{x\to0}\frac{f(x)-f(0)}x
=\lim_{x\to0}\frac{f(x)}x=L,
\]
所以 B 正确。

由 \(f(x)/x\) 有极限可知该商在 0 的去心邻域内有界，因此存在 \(C>0\)，使
\[
|f(x)|<C|x|,
\]
D 正确。

C 错误。例如
\[
f(x)=x
\]
满足题设，但
\[
f'(x)=1,
\]
所以
\[
\lim_{x\to0}f'(x)=1\ne0.
\]
\answer{C}
\examnote{函数在一点可导，不代表去心邻域内的导函数极限等于该点导数，更不代表等于 0。}
\end{solutionblock}

\section{凸性与级数保号估计}

\begin{problemblock}
\textbf{例17.} 设函数 \(f(x)\) 在 \(x=x_0\) 的某邻域 \(U\) 内存在连续的二阶导数。

\begin{enumerate}[label=(\arabic*)]
\item 设当 \(h>0,\ x_0-h\in U,\ x_0+h\in U\) 时，恒有
\[
f(x_0)<\frac12\bigl[f(x_0+h)+f(x_0-h)\bigr],\tag{*}
\]
证明 \(f''(x_0)\ge0\)。
\item 如果 \(f''(x_0)>0\)，证明必存在 \(h>0,\ x_0-h\in U,\ x_0+h\in U\)，使 \((*)\) 式成立。
\end{enumerate}
\end{problemblock}
\begin{solutionblock}
\analysis{这是用中心二阶差分刻画二阶导数。中心平均大于中点函数值，说明二阶差分非负。}
(1) 由 \((*)\) 得
\[
f(x_0+h)-2f(x_0)+f(x_0-h)>0.
\]
两边除以 \(h^2>0\)，有
\[
\frac{f(x_0+h)-2f(x_0)+f(x_0-h)}{h^2}>0.
\]
令 \(h\to0^+\)。由于 \(f\) 二阶可导，
\[
\lim_{h\to0}
\frac{f(x_0+h)-2f(x_0)+f(x_0-h)}{h^2}
=f''(x_0).
\]
因此
\[
f''(x_0)\ge0.
\]

(2) 若 \(f''(x_0)>0\)，由 \(f''\) 连续，存在 \(\delta>0\)，使得当
\[
|x-x_0|<\delta
\]
时
\[
f''(x)>0.
\]
取 \(0<h<\delta\)，且 \(x_0\pm h\in U\)。由 Taylor 公式，
\[
f(x_0+h)=f(x_0)+f'(x_0)h+\frac12f''(\xi_1)h^2,
\]
\[
f(x_0-h)=f(x_0)-f'(x_0)h+\frac12f''(\xi_2)h^2,
\]
其中 \(\xi_1\in(x_0,x_0+h)\)，\(\xi_2\in(x_0-h,x_0)\)。
两式相加得
\[
f(x_0+h)+f(x_0-h)-2f(x_0)
=\frac12\bigl[f''(\xi_1)+f''(\xi_2)\bigr]h^2>0.
\]
故
\[
f(x_0)<\frac12\bigl[f(x_0+h)+f(x_0-h)\bigr],
\]
即 \((*)\) 成立。
\answer{(1) \(f''(x_0)\ge0\)；(2) 存在这样的 \(h>0\)}
\examnote{中心二阶差分
\(\bigl[f(x_0+h)-2f(x_0)+f(x_0-h)\bigr]/h^2\)
是判断凸性的常用工具。}
\end{solutionblock}

\begin{problemblock}
\textbf{例18.} 若 \(\{a_n\}\) 收敛，\(\displaystyle \sum_{n=1}^{\infty}b_n\) 绝对收敛，证明：
\[
\sum_{n=1}^{\infty}a_nb_n
\]
绝对收敛。
\end{problemblock}
\begin{solutionblock}
\analysis{收敛数列必有界。把 \(a_n\) 的绝对值用一个常数控制，就能与 \(\sum |b_n|\) 比较。}
因为 \(\{a_n\}\) 收敛，所以 \(\{a_n\}\) 有界。存在 \(M>0\)，使得
\[
|a_n|\le M,\qquad n=1,2,\cdots.
\]
又因为
\[
\sum_{n=1}^{\infty}b_n
\]
绝对收敛，所以
\[
\sum_{n=1}^{\infty}|b_n|<+\infty.
\]
于是
\[
\sum_{n=1}^{\infty}|a_nb_n|
\le
\sum_{n=1}^{\infty}M|b_n|
=M\sum_{n=1}^{\infty}|b_n|
<+\infty.
\]
因此
\[
\sum_{n=1}^{\infty}a_nb_n
\]
绝对收敛。
\answer{已证 \(\displaystyle \sum_{n=1}^{\infty}a_nb_n\) 绝对收敛}
\examnote{“收敛数列有界 + 绝对收敛级数可比较”是级数证明题里最常用的一步。}
\end{solutionblock}
